\chapter{Kết luận và Hướng Phát triển}
\label{ch:conclusion}

\section{Kết quả hiện thực}
\label{sec:conclusion_summary}

Đề tài đã hiện thực \textbf{AST-Anchored Neuro-Symbolic Verification Framework} cho một tập con Python. LLM đề xuất khẳng định và bất biến ứng viên. Hệ thống xác định vị trí kiểm tra từ AST, dựng quan hệ trạng thái bằng thực thi ký hiệu và dùng Z3 để xét tính khả đạt và tìm vi phạm.

Trên 200 chương trình vòng lặp được đối chiếu bằng khẳng định gốc, Proposed với MiniCPM và Proposed Offline đều đạt 200/200 phán quyết đúng và tái hiện được 100/100 phản ví dụ trong mỗi chế độ. Kết quả đồng điểm giữa hai chế độ Proposed cho thấy độ chính xác trên tập này chưa phân biệt được tác động riêng của MiniCPM.

\section{Đóng góp của đề tài}
\label{sec:contributions_highlight}

\begin{enumerate}
    \item \textbf{Neo AST và giao thức hai truy vấn:} Hệ thống gắn claim với câu lệnh đích và kiểm tra tính khả đạt trước khi tìm vi phạm trên đường đã mô hình hóa.
    \item \textbf{Replay và phản hồi hồi cứu:} Hệ thống kiểm tra phản ví dụ bằng CPython và có thể dùng trạng thái quan sát được để yêu cầu một claim khác, sau đó kiểm tra lại claim đó.
    \item \textbf{Hạ tầng đánh giá:} Đề tài chuẩn bị bộ dữ liệu \texttt{bench\_v1}, tập khẳng định dùng chung và 14 baseline. Nhóm công bố bảng kết quả trên tập 200 mẫu.
\end{enumerate}

\section{Hạn chế của Đề tài}
\label{sec:limitations}

Các hạn chế kỹ thuật của hệ thống gồm những điểm sau.
\begin{itemize}
    \item \textbf{Phạm vi Ngữ nghĩa Ngôn ngữ Hẹp:} Hệ thống hiện chỉ hỗ trợ tập con QF-LIA của Python (số nguyên toán học, biểu thức tuyến tính). Các tính năng hiện đại và phổ biến của Python như lập trình hướng đối tượng, cấu trúc dữ liệu động (danh sách lồng nhau, từ điển), số thực dấu phẩy động, xử lý chuỗi ký tự, và đệ quy sâu vẫn chưa được mô hình hóa trong TCB.
    \item \textbf{Bộ Chuyển đổi AST-to-SMT Chưa Được Chứng minh Cơ giới hóa:} Bộ chuyển đổi có kiểm thử hồi quy, nhưng tính tương ứng giữa phép dịch và ngữ nghĩa Python chưa được chứng minh bằng một trợ lý chứng minh định lý.
    \item \textbf{Thách thức trong Tổng hợp Bất biến Phức tạp:} Mặc dù đã tích hợp $k$-induction và đề xuất bất biến từ LLM, đối với các vòng lặp chứa các điều kiện rẽ nhánh phi tuyến hoặc các mối quan hệ biến đan xen phức tạp, solver vẫn có thể rơi vào tình trạng timeout hoặc trả về \texttt{UNKNOWN\_TIMEOUT}.
\end{itemize}

Các cấu hình baseline được chọn theo kết quả trên cùng bộ 200 mẫu và ngân sách thực thi khác nhau. Chương \ref{ch:evaluation} nêu phạm vi của các kết luận và yêu cầu đối với phép thử trên dữ liệu giữ riêng.

\section{Hướng Phát triển Tương lai}
\label{sec:future_work}

Từ những kết quả và hạn chế đã được nhận diện, nhóm nghiên cứu đề xuất các định hướng phát triển tiếp theo trong tương lai:

\begin{enumerate}
    \item \textbf{Mở rộng Ngữ nghĩa Sang Lý thuyết Mảng và Cấu trúc Dữ liệu Động:} Tích hợp lý thuyết mảng không lượng từ (Quantifier-Free Array Theory -- QF-AX) trong SMT-LIB2 vào Symbolic Executor để hỗ trợ đầy đủ các thao tác đọc ghi mảng và danh sách Python (\texttt{list[int]}), mở rộng khả năng kiểm chứng cho các thuật toán sắp xếp và tìm kiếm cổ điển.
    \item \textbf{Tối ưu hóa Hiệu năng và Giảm Độ trễ:}
    \begin{itemize}
        \item Áp dụng kỹ thuật tỉa nhánh không khả đạt sớm (Early Dead-code Pruning): Khi một guard nhánh được Z3 xác định là \texttt{unsat}, toàn bộ cây con cú pháp bên dưới sẽ bị loại bỏ ngay lập tức khỏi quy trình phân tích SSA.
        \item Caching ngữ nghĩa SMT: Lưu trữ băm (hash) của các mệnh đề đường đi tương đương để tái sử dụng kết quả giải của Z3 giữa các lần chạy liên tiếp.
        \item Thu gọn Ngữ cảnh Phản hồi (Feedback Context Compaction): Rút trích chỉ các biến liên quan trực tiếp đến lát cắt lỗi (failure slice) thay vì gửi toàn bộ \texttt{locals()}, giúp giảm lượng token tiêu thụ và tăng tốc độ xử lý của LLM.
    \end{itemize}
    \item \textbf{Phân tích Tính Dừng và Hàm Giảm (Termination Analysis):} Bổ sung mô-đun yêu cầu LLM đề xuất các hàm giảm (Ranking Functions / Variants) $\mathcal{R}: S \rightarrow \mathbb{N}$ và sử dụng solver để chứng minh $\mathcal{R}(s') < \mathcal{R}(s)$, nâng cấp phán quyết từ tính đúng đắn từng phần (Partial Correctness) lên tính đúng đắn toàn phần (Total Correctness).
    \item \textbf{Cơ giới hóa và Kiểm chứng Bản thân TCB trong Lean 4:} Tiến hành hình thức hóa toàn bộ ngữ nghĩa của Safe Subset và chứng minh tính đúng đắn của Trình thực thi Ký hiệu (Symbolic Executor Soundness Proof) bên trong trợ lý chứng minh định lý Lean 4, đạt tới cấp độ tin cậy cao nhất trong khoa học máy tính.
    \item \textbf{Ứng dụng Thực tiễn trong Giáo dục và Công nghiệp:} Đóng gói framework thành tiện ích mở rộng cho Visual Studio Code (VS Code Extension) hỗ trợ lập trình viên phát hiện lỗi logic ngay trong quá trình soạn thảo, đồng thời ứng dụng làm hệ thống chấm bài và hướng dẫn sửa bài tập lập trình tự động trong các khóa học Khoa học Máy tính tại các trường đại học.
\end{enumerate}
